\editionPage[\pdfbookmark[0]{Deux bilans de fin d’année}{bm-bilans}]{bilans-guide}{Travailler les bilans}{Seconde · Synthèse des apprentissages}{Deux entraînements originaux, sans valeur d’épreuve officielle.}
\begin{blocEdition}{bilans-guide}
\editionEtape{Organisation proposée}
Chaque bilan dure environ 2 heures. Les dix automatismes se traitent sans calculatrice en 15 à 20 minutes. Pour les problèmes, la calculatrice est autorisée si le professeur ne donne pas d’autre consigne. Les questions demandant une forme exacte restent à traiter exactement.
\editionEtape{Barème d’entraînement}
Les automatismes valent 5 points, à raison de 0,5 point chacun. Les trois problèmes valent 5 points chacun. Total : 20 points ; le barème détaillé est un choix de ce livre. Une démarche correcte peut être valorisée même si un calcul intermédiaire échoue.
\editionEtape{Avant et après la séance}
Prévoir une feuille de brouillon et une copie pour la rédaction. Donner les hypothèses, les étapes et les unités. Après correction, classer les difficultés : notion, méthode, calcul ou rédaction. Refaire une question de chaque difficulté repérée.
\editionRetenir{Un bilan sert à identifier le prochain travail utile. Reprendre ensuite les exercices ciblés.}
\end{blocEdition}
\editionPage{bilan-A-auto}{Bilan A · Automatismes}{Fonctions, géométrie et conditionnement}{5 points · Sans calculatrice · 0,5 point par réponse}
\begin{blocEdition}{bilan-A-auto}
\begin{questions}
\item Calculer $3/4-1/6$.
\item Simplifier $\sqrt{75}$.
\item Développer $(x-4)^2$.
\item Résoudre $-2x\le6$.
\item Écrire $|x-1|\le2$ sous forme d’intervalle.
\item Calculer l’image de $-2$ par $f(x)=3x+1$.
\item Donner le domaine de $1/(x+5)$.
\item Donner le milieu de $A(0;2)$ et $B(6;4)$.
\item Une hausse de 10\,\% a quel coefficient multiplicateur ?
\item Pour $P(A)=0{,}35$, calculer $P(\overline A)$.
\end{questions}
\end{blocEdition}
\editionPage{bilan-A-ex-1}{Bilan A · Problème 1}{L’aire d’un rectangle}{5 points · Démarche et rédaction attendues}
\begin{blocEdition}{bilan-A-ex-1}
On étudie $f(x)=x(6-x)$ sur $[0;6]$. Pour $0<x<6$, cette fonction est l’aire d’un rectangle de côtés $x$ et $6-x$, en centimètres.\begin{questions}
\item (1 point) Calculer $f(0)$, $f(2)$ et $f(6)$.
\item (1 point) Montrer $f(x)=9-(x-3)^2$.
\item (1 point) Donner les variations et le maximum sur $[0;6]$.
\item (1 point) Résoudre $f(x)\ge8$.
\item (1 point) Conclure sur les dimensions et l’aire maximale du rectangle ; préciser son périmètre.
\end{questions}\label{fig-bilan-A-ex-1}
\figureGeo{bilan-A-aire}

\end{blocEdition}
\editionPage{bilan-A-ex-2}{Bilan A · Problème 2}{Deux médianes se rencontrent}{5 points · Démarche et rédaction attendues}
\begin{blocEdition}{bilan-A-ex-2}
Dans un repère orthonormé, $A(0;0)$, $B(8;0)$ et $C(2;6)$. On note $I$ le milieu de $[BC]$ et $J$ celui de $[AC]$.\begin{questions}
\item (1 point) Calculer $I$ et $J$.
\item (1 point) Donner une équation de $(AI)$.
\item (1 point) Donner une équation de $(BJ)$.
\item (1 point) Déterminer leur intersection $G$.
\item (1 point) Montrer que $\vect{AG}=\dfrac23\vect{AI}$ et expliquer ce que cette égalité apporte.
\end{questions}\label{fig-bilan-A-ex-2}
\figureGeo{bilan-A-geometrie}

\end{blocEdition}
\editionPage{bilan-A-ex-3}{Bilan A · Problème 3}{Une enquête et son arbre}{5 points · Démarche et rédaction attendues}
\begin{blocEdition}{bilan-A-ex-3}
Une enquête construite pour ce bilan porte sur 200 personnes. Le groupe $A$ en compte 80, dont 48 utilisent le service $S$. Le groupe $B$ en compte 120, dont 36 utilisent $S$. On choisit uniformément une personne.\begin{questions}
\item (1 point) Dresser le tableau complet $A,B$ / $S,\overline S$.
\item (1 point) Calculer la fréquence globale de $S$.
\item (1 point) Calculer $P_A(S)$ et $P_B(S)$.
\item (1 point) Compléter l’arbre et calculer le chemin $A$ puis $S$.
\item (1 point) Calculer la part de $A$ parmi les utilisateurs et expliquer pourquoi elle ne vaut pas $P_A(S)$.
\end{questions}\label{fig-bilan-A-ex-3}
\figureGeo{bilan-A-arbre}

\end{blocEdition}
\editionPage{bilan-B-auto}{Bilan B · Automatismes}{Statistiques, configurations et tirages}{5 points · Sans calculatrice · 0,5 point par réponse}
\begin{blocEdition}{bilan-B-auto}
\begin{questions}
\item Rendre $42/63$ irréductible.
\item Calculer $2^3\times2^{-5}$.
\item Factoriser $x^2-9$.
\item Résoudre $(x-2)(x+1)=0$.
\item Résoudre $\sqrt x=3$.
\item Donner la pente de $y=-3x+4$.
\item Calculer $\det((2;1);(4;3))$.
\item Une hausse de 25\,\% suivie d’une baisse de 20\,\% donne quel taux global ?
\item Donner la médiane de $1,3,5,9,12$.
\item Combien de valeurs parcourt \code{range(2,7)} ?
\end{questions}
\end{blocEdition}
\editionPage{bilan-B-ex-1}{Bilan B · Problème 1}{Des durées regroupées}{5 points · Démarche et rédaction attendues}
\begin{blocEdition}{bilan-B-ex-1}
Des durées en minutes sont regroupées en classes.\begin{center}\renewcommand{\arraystretch}{1.35}\begin{tabular}{c|ccc}\toprule
Classe & $[0;10[$ & $[10;20[$ & $[20;30]$\\\midrule
Effectif & 4 & 8 & 8\\
\bottomrule\end{tabular}\end{center}\begin{questions}
\item (1 point) Donner l’effectif total et les fréquences cumulées aux bornes 10, 20 et 30.
\item (1 point) Estimer la moyenne par les centres, en précisant l’hypothèse.
\item (1 point) Déterminer la classe médiane et estimer la médiane sous l’hypothèse uniforme.
\item (1 point) Les moyennes réelles de classe sont 6, 14 et 27. Calculer la moyenne exacte.
\item (1 point) On ajoute à cette dernière série une durée de 60 minutes. Calculer la nouvelle moyenne et commenter.
\end{questions}\label{fig-bilan-B-ex-1}
\figureGeo{bilan-B-stats}

\end{blocEdition}
\editionPage{bilan-B-ex-2}{Bilan B · Problème 2}{Un rectangle et son cercle}{5 points · Démarche et rédaction attendues}
\begin{blocEdition}{bilan-B-ex-2}
On donne $A(1;1)$, $B(5;1)$, $C(1;4)$ et $D(5;4)$ dans un repère orthonormé.\begin{questions}
\item (1 point) Montrer que le quadrilatère $ABDC$ est un rectangle et calculer son aire.
\item (1 point) Calculer le milieu $O$ de $[BC]$.
\item (1 point) Calculer $OA$, $OB$, $OC$ et $OD$.
\item (1 point) Donner le centre et le rayon du cercle passant par les quatre points.
\item (1 point) Donner des équations des deux médiatrices de $[AB]$ et $[AC]$ et vérifier leur intersection.
\end{questions}\label{fig-bilan-B-ex-2}
\figureGeo{bilan-B-geometrie}

\end{blocEdition}
\editionPage{bilan-B-ex-3}{Bilan B · Problème 3}{Une urne sans remise}{5 points · Démarche et rédaction attendues}
\begin{blocEdition}{bilan-B-ex-3}
Une urne contient quatre boules rouges et une bleue. On tire deux boules successivement, sans remise, chaque boule restante ayant la même chance d’être choisie. On note $R_1$ et $B_1$ les couleurs du premier tirage, $R_2$ et $B_2$ celles du second.\begin{questions}
\item (1 point) Donner les probabilités du premier tirage.
\item (1 point) Compléter les branches du second tirage, y compris le chemin impossible.
\item (1 point) Calculer $P(R_1\cap R_2)$.
\item (1 point) Calculer $P(R_1\cap B_2)$ et $P(B_1\cap R_2)$.
\item (1 point) On sait que la seconde boule est bleue. Que peut-on affirmer de la première ? Donner la probabilité conditionnelle correspondante.
\end{questions}\label{fig-bilan-B-ex-3}
\figureGeo{bilan-B-arbre}

\end{blocEdition}
